\documentclass{llncs}

\usepackage{amsmath,amssymb}
\usepackage{graphicx}
\usepackage[hidelinks]{hyperref}

\begin{document}

\title{Paper Title Goes Here}
\titlerunning{Short Paper Title}

\author{First Author\inst{1} \and Second Author\inst{2}}
\authorrunning{F. Author and S. Author}

\institute{
    Department of Computer Science, University Name, City, Country \\
    \email{first.author@example.com} \and
    Department of Computer Science, Other University, City, Country \\
    \email{second.author@example.com}
}

\maketitle

\begin{abstract}
This document is a starting point for a Springer LNCS conference submission. Replace this abstract with a concise summary (150--250 words) of the problem, approach, and key results of the work.
\keywords{keyword one \and keyword two \and keyword three}
\end{abstract}

\section{Introduction}
Introduce the problem, motivate why it matters, summarize prior approaches and their limitations, and state the paper's contributions as a short itemized list.

\section{Related Work}
Summarize and position the work relative to the most relevant prior literature, organized by theme rather than chronologically.

\section{Proposed Method}
Describe the proposed method, model, or system.

\subsection{Problem Formulation}
State assumptions and formal notation here.

\subsection{Approach}
Describe the core technical contribution here.

\section{Experiments}
Describe the experimental setup: datasets, baselines, evaluation metrics, and implementation details needed for reproducibility.

\section{Conclusion}
Summarize findings, discuss limitations, and outline directions for future work.

\begin{thebibliography}{8}

\bibitem{ref1}
Author, A., Author, B.: Title of the referenced paper. Journal Name 1(1), 1--10 (Year)

\bibitem{ref2}
Author, C., Author, D., Author, E.: Title of another referenced paper. In: Proc. Conference Name, pp. 1--6 (Year)

\end{thebibliography}

\end{document}
